% ============================================================
\section{Formative Interview Plan}
\label{app:formative-interview-plan}

This section documents the semi-structured interview plan the six formative sessions followed. Each session combined questions about current study-planning practice with a think-aloud walkthrough of the prototype, in order to understand how students plan and to gather feedback on the hybrid graph-table concept. The sessions were conducted in German; the script below is the author's translation.

\subsection{Overview}

\begin{table}[!t]
\centering
\begin{tabular}{p{4.2cm}p{2.2cm}p{7.2cm}}
\hline
\textbf{Phase} & \textbf{Duration} & \textbf{Goal} \\
\hline
A. Welcome \& Consent & 3 min & Build trust, explain study purpose \\
B. Current Planning Behaviour & 15 min & Understand existing workflows and pain points \\
C. Prototype Exploration (Think-Aloud) & 30 min & Observe expectations, reactions, and design implications \\
D. Reflection \& Wrap-Up & 12 min & Summarise impressions and collect improvement ideas \\
\hline
\end{tabular}
\caption{Structure of the interview sessions. The four phases were held to in every session; the durations are the planned budget rather than a measured mean, and the script for each phase is reproduced in Section~\ref{app:formative-script}.}
\label{tab:interview-structure}
\end{table}

The study is the formative one reported in Chapter~\ref{chap:formative-study}: six participants, three from the Bachelor's and three from the Master's programme, in a semi-structured interview with a think-aloud prototype walkthrough, recorded as audio. Each session ran to the four phases of Table~\ref{tab:interview-structure}, and Section~\ref{app:formative-script} reproduces the script phase by phase.

\subsection{Detailed Interview Script}
\label{app:formative-script}

% ------------------------------------------------------------
\subsubsection{Phase A: Welcome, Framing, and Consent}

\paragraph{Introductory script.}
\begin{quote}
Thanks for taking part. I am testing an early concept for a study-planning tool that combines a Graph View and a Table View with recommendations.

The topic is study planning. What do you think this means? Could you explain it in your own words so that we know we are talking about the same thing?
\end{quote}

\paragraph{Clarification of the study scope.}
\begin{quote}
In our case, we are interested in how you plan your whole study programme. That means we focus on the course level:
\begin{itemize}
    \item how your courses are connected,
    \item which dependencies or prerequisites exist,
    \item which courses you want to take, should take, or have to take,
    \item and in which semester you take them.
\end{itemize}
We are not looking at how you plan your time or resources within a single course. Instead, we care about this higher-level, programme-wide planning.
\end{quote}

\paragraph{Consent and background questions.}
\begin{quote}
There are no right or wrong answers; we are interested in your honest opinions. With your permission, I will record audio for analysis; all data will be anonymous.
\end{quote}

\noindent
\textbf{Background questions:}
\begin{itemize}
    \item Which programme and semester are you in?
    \item How many ECTS are you still missing approximately?
    \item Do you agree to audio recording?
\end{itemize}

% ------------------------------------------------------------
\subsubsection{Phase B: Current Planning Behaviour}

\paragraph{Core prompts.}
\begin{itemize}
    \item Walk me through the last time you planned your semester. What was the first thing you did?
    \item What kind of information did you look for, or need, in order to decide?
    \item What was most important to you while planning your semester?
    \item Which tools or resources did you use for planning? Can you show or describe one of them?
    \item What does each tool help you with specifically? What do you mainly use it for: overview, scheduling, checking rules, workload balance, or something else?
    \item Where did you run into challenges using these tools? Can you give an example of when something did not work as expected, and how you handled it?
    \item Have you ever been frustrated about parts of the process?
    \item What is missing from your current tools or process? Which tasks do you still have to do manually, or outside of those tools?
\end{itemize}


\paragraph{Recommendation-related prompts.}
\begin{itemize}
    \item If a tool could give you recommendations while you plan, what kind of recommendations would be interesting?
    \item What should such recommendations look like?
    \item When should they appear: during course selection, after choosing, or somewhere else?
\end{itemize}

% ------------------------------------------------------------
\subsubsection{Phase C: Prototype Exploration (Think-Aloud)}

\paragraph{Method instruction.}
\begin{quote}
Now please look at this prototype with me. Imagine you are planning next semester. Please say aloud what you expect, what you like, and what confuses you.
\end{quote}

The walkthrough ran in six demonstration steps, three in the Table View and then three in the Graph View, and the moderator paused after every step so that participants could react to what they had just seen. Table~\ref{tab:walkthrough-steps} lists what was shown at each step. Two points carried a question beyond the free comment the think-aloud invited: after the graph was first shown, participants were asked whether they could think of other ways to connect the graph with the table, and the walkthrough closed on the four questions below.


\begin{table}[htbp]
    \centering
    \small
    \begin{tabular}{@{}c>{\raggedright\arraybackslash}p{12.6cm}@{}}
        \toprule
        Step & Shown \\
        \midrule
        \multicolumn{2}{@{}l}{\emph{Table View}} \\
        1 & The overview; the search field; exam subjects in their distinct colours; module numbers within an exam subject; the semester grid with its per-semester ECTS sum \\
        2 & The sidebar, expanded; the difference between exam-subject, module and course items; the M, C and E markers; what the tool tracks \\
        3 & Drag interactions; how the ECTS sum responds to a change; how a module is moved \\
        \midrule
        \multicolumn{2}{@{}l}{\emph{Graph View}} \\
        4 & The graph at programme level, including the add-to-plan interaction; the stars; the M, C and E markers; the filter controls; exam-subject, module and course items \\
        5 & The filter controls in use: exam subjects, only added courses, content similarities, formal prerequisites, show all, collapse all \\
        6 & Similar courses \\
        \bottomrule
    \end{tabular}
    \caption{The prototype walkthrough of Phase~C, in the order it was demonstrated. Steps~1 to~3 were shown in the Table View and steps~4 to~6 in the Graph View.}
    \label{tab:walkthrough-steps}
\end{table}

\paragraph{Closing discussion: recommendations.}
The walkthrough ended on the recommendation concept, which the prototype sketched rather than implemented. Four questions were put:
\begin{itemize}
    \item On what could recommendations be based?
    \item How could recommendations be integrated in the graph, and how in the table?
    \item What would you need in order to trust them?
\end{itemize}

% ------------------------------------------------------------
\subsubsection{Phase D: Reflection and Wrap-Up}

\begin{itemize}
    \item If you could add one thing to the tool, what would it be?
    \item Would you integrate a tool with the shown concepts into your personal study planning, and why?
    \item Follow-up (if not): What would increase the chance that you would do so?
    \item Is there anything else you would like to share with us?
    \item We would contact you for a follow-up if we build the actual tool. Would this be okay?
\end{itemize}

% ============================================================